\documentclass[14pt,a4paper,oneside]{extbook}
\usepackage[utf8]{vietnam}
\usepackage{amsmath,amssymb,mathrsfs,fancyhdr,enumerate,multirow,makecell}
\usepackage{tkz-euclide,tikz-3dplot,tikz,tkz-tab}
\usetikzlibrary{arrows.meta,calc,intersections,angles,quotes,backgrounds,shapes.geometric,patterns,through,shadows,positioning}
\usepackage{yhmath}
\usepackage[top=0.7cm, bottom=0.7cm, left=1.0cm, right=1.0cm, headheight=16pt, headsep=3mm, footskip=5mm]{geometry}
\usepackage[hidelinks,unicode]{hyperref}
\usepackage{qrcode}
\usepackage[many]{tcolorbox}
\usepackage{fontawesome}

\pagestyle{fancy}
\fancyhf{}
\fancyhead[L]{\footnotesize\bfseries\color{teal!85!black}\faBook~SỔ TAY KIẾN THỨC TOÁN 4}
\fancyhead[R]{\footnotesize\bfseries\color{blue!85!black}Lớp Toán Thầy Lĩnh -- loptoanthaylinh.vn}
\fancyfoot[C]{\footnotesize\bfseries\color{gray!80!black}-- Sổ tay kiến thức: Bài 65. Tìm phân số của một số --}
\renewcommand{\headrulewidth}{0.4pt}

\begin{document}

% --- BANNER TIÊU ĐỀ BÀI HỌC ---
\noindent
\begin{tcolorbox}[
  enhanced,
  colback=teal!8!white,
  colframe=teal!80!black,
  arc=3mm, boxrule=1pt,
  left=3mm, right=3mm, top=2mm, bottom=2mm,
  halign=center
]
  {\Large\bfseries\color{teal!90!black}\faBookmark\ BÀI 65: TÌM PHÂN SỐ CỦA MỘT SỐ}\\[1.5pt]
  {\footnotesize\itshape\color{black!75} Bản chất chia phần bằng nhau, quy tắc nhân số với phân số và ứng dụng thực tế trên 1 trang A4}
\end{tcolorbox}

\vspace{-2.5mm}

\noindent
\begin{minipage}[t]{0.49\textwidth}\vspace{0pt}
\begin{tcolorbox}[
  enhanced,
  colback=blue!4!white,
  colframe=blue!75!black,
  arc=2mm, boxrule=0.8pt,
  left=2mm, right=2mm, top=1.8mm, bottom=1.8mm,
  title={\footnotesize\bfseries\color{white}\faCalculator\ 1. Quy tắc cốt lõi}
]
{\footnotesize
Muốn tìm $\mathbf{\dfrac{m}{n}}$ của số $A$, ta \textbf{lấy số $A$ nhân với phân số $\dfrac{m}{n}$}:
$$\text{Giá trị cần tìm} = A \times \dfrac{m}{n} = \dfrac{A \times m}{n}$$
\textbf{Cách nhẩm bản chất 2 bước:}
\begin{itemize}\setlength{\itemsep}{1pt}\setlength{\leftskip}{-2.5mm}
  \item Bước 1: Chia $A$ cho mẫu số $n$ để tìm \textbf{giá trị 1 phần}.
  \item Bước 2: Lấy giá trị 1 phần nhân với tử số $m$ để tìm \textbf{giá trị $m$ phần}.
\end{itemize}
}
\end{tcolorbox}
\end{minipage}
\hfill
\begin{minipage}[t]{0.49\textwidth}\vspace{0pt}
\begin{tcolorbox}[
  enhanced,
  colback=teal!4!white,
  colframe=teal!75!black,
  arc=2mm, boxrule=0.8pt,
  left=2mm, right=2mm, top=1.8mm, bottom=1.8mm,
  title={\footnotesize\bfseries\color{white}\faTags\ 2. Các dạng toán điển hình}
]
{\footnotesize
\begin{itemize}\setlength{\itemsep}{1.5pt}\setlength{\leftskip}{-2.5mm}
  \item \textbf{Dạng 1 (Đo lường, khối lượng, diện tích):}\\
  Tìm $\dfrac{2}{5}$ của $30\text{ kg} \implies 30 \times \dfrac{2}{5} = 12\text{ kg}$.
  \item \textbf{Dạng 2 (Thời gian, tiền bạc):}\\
  $\dfrac{3}{4}$ của 1 giờ ($=60$ phút) $\implies 60 \times \dfrac{3}{4} = 45$ phút.
  \item \textbf{Dạng 3 (Số lượng học sinh, đồ vật):}\\
  Số học sinh nữ bằng $\dfrac{3}{5}$ số học sinh cả lớp...
\end{itemize}
}
\end{tcolorbox}
\end{minipage}

\vspace{-2.5mm}
% --- KHỐI 3: SƠ ĐỒ TƯ DUY TIKZ ---
\noindent
\begin{tcolorbox}[
  enhanced,
  colback=orange!4!white,
  colframe=orange!80!black,
  arc=2mm, boxrule=0.8pt,
  left=2mm, right=2mm, top=1.5mm, bottom=1.5mm,
  title={\footnotesize\bfseries\color{white}\faSitemap\ SƠ ĐỒ TƯ DUY: MÔ HÌNH THANH ĐOẠN THẲNG TÌM $\frac{3}{5}$ CỦA 25}
]
\centering
\begin{tikzpicture}[scale=0.9, transform shape]
  % Toàn bộ thanh đoạn thẳng là 25
  \draw[thick, fill=gray!10] (0, 0.5) rectangle (10, 1.3);
  \foreach \x in {2, 4, 6, 8} {
    \draw[thick] (\x, 0.5) -- (\x, 1.3);
  }
  % Tô màu 3 phần (3/5)
  \fill[teal!40] (0, 0.5) rectangle (6, 1.3);
  \foreach \x in {2, 4} {
    \draw[thick] (\x, 0.5) -- (\x, 1.3);
  }

  % Kích thước tổng
  \draw[decorate, decoration={brace, amplitude=4pt, mirror}, thick] (0, 0.3) -- (10, 0.3) node[midway, below=3pt, font=\scriptsize\bfseries] {Cả đoạn thẳng = 25 (chia 5 phần bằng nhau)};

  % Phần cần tìm (3 phần)
  \draw[decorate, decoration={brace, amplitude=4pt}, thick, teal!90!black] (0, 1.5) -- (6, 1.5) node[midway, above=3pt, font=\scriptsize\bfseries\color{teal!90!black}] {Cần tìm $\dfrac{3}{5}$ của $25 = 25 \times \dfrac{3}{5} = \mathbf{15}$};

  % Giá trị 1 phần
  \node[font=\tiny\bfseries\color{gray!80!black}] at (1, 0.9) {5};
  \node[font=\tiny\bfseries\color{gray!80!black}] at (3, 0.9) {5};
  \node[font=\tiny\bfseries\color{gray!80!black}] at (5, 0.9) {5};
  \node[font=\tiny\bfseries\color{gray!80!black}] at (7, 0.9) {5};
  \node[font=\tiny\bfseries\color{gray!80!black}] at (9, 0.9) {5};
\end{tikzpicture}
\end{tcolorbox}

\vspace{-2.5mm}
% --- KHỐI 4: VÍ DỤ & MẸO TRÁNH BẪY ---
\noindent
\begin{minipage}[t]{0.49\textwidth}\vspace{0pt}
\begin{tcolorbox}[
  enhanced,
  colback=green!4!white,
  colframe=green!60!black,
  arc=2mm, boxrule=0.8pt,
  left=2mm, right=2mm, top=1.8mm, bottom=1.8mm,
  title={\footnotesize\bfseries\color{white}\faCheckCircle\ Ví dụ thực hành chuẩn mực}
]
{\footnotesize
\textbf{Bài toán:} Một lớp có 35 học sinh, trong đó $\dfrac{3}{5}$ số học sinh là nữ. Hỏi lớp đó có bao nhiêu học sinh nữ?\\
\textbf{Bài giải:}\\
Số học sinh nữ của lớp đó là:
$$35 \times \dfrac{3}{5} = \mathbf{21} \text{ (học sinh)}$$
\textit{Hoặc tính nhẩm:} $35 : 5 \times 3 = 21$ (học sinh).\\
\textbf{Đáp số:} $21$ học sinh nữ.
}
\end{tcolorbox}
\end{minipage}
\hfill
\begin{minipage}[t]{0.49\textwidth}\vspace{0pt}
\begin{tcolorbox}[
  enhanced,
  colback=red!4!white,
  colframe=red!70!black,
  arc=2mm, boxrule=0.8pt,
  left=2mm, right=2mm, top=1.8mm, bottom=1.8mm,
  title={\footnotesize\bfseries\color{white}\faExclamationTriangle\ Mẹo tránh bẫy trắc nghiệm}
]
{\footnotesize
\begin{itemize}\setlength{\itemsep}{1.5pt}\setlength{\leftskip}{-2.5mm}
  \item \textbf{Bẫy nhầm phép tính nhân thành chia:}\\
  Tìm phân số của một số là \textbf{phép nhân} ($A \times \frac{m}{n}$), nhiều bạn hay làm nhầm thành phép chia ($A : \frac{m}{n}$).
  \item \textbf{Bẫy chưa đổi đơn vị:}\\
  Nếu đề cho "Tìm $\frac{2}{3}$ của 2 giờ theo phút", phải đổi $2\text{ giờ} = 120\text{ phút}$ rồi mới nhân: $120 \times \frac{2}{3} = 80\text{ phút}$.
\end{itemize}
}
\end{tcolorbox}
\end{minipage}

\end{document}
